\logosection{\faUser}{个人简介}
东华大学软件工程 2027 届，求职方向为 AI 应用工程与 Java 后端。主导 408Master 教育 AI 平台和 StoryCraft AI 文字冒险游戏，具备从需求分析、模型接入、RAG、对话记忆、Tool Calling 到全栈交付的实践。重视模型能力与确定性业务代码的边界，能够结合成本、并发、稳定性和安全解释技术选型，而不把简单 API 调用包装成 Agent。

\logosection{\faGraduationCap}{教育经历}
\datedline{\textbf{东华大学（211）}}{\dateRange{2023.09}{2027.06}}
\datedline{\tripleInfo{软件工程}{本科 · 2027 届}{计算机科学与技术学院}}{上海}
\datedline{主修课程：数据结构、算法、操作系统、计算机网络、数据库、软件工程、自然语言处理}{}

\logosection{\faCogs}{核心技能}
\begin{itemize}[parsep=0.5ex]
  \item \textbf{AI 应用}：Spring AI 2（ChatModel、ChatClient、Advisor、ChatMemory、VectorStore、Tool Calling），理解 Prompt、上下文窗口、Embedding、RAG、工具执行循环及人工确认边界。
  \item \textbf{Java 后端}：Java 21、Spring Boot 4、Spring MVC、Spring Security、MyBatis、MySQL、Redis、REST、SSE、Maven；能够进行旧系统渐进迁移、接口兼容和事务边界设计。
  \item \textbf{RAG 与数据}：Qdrant、语义分块、overlap、Top-K、相似度检索、metadata 与引用溯源；理解 MySQL 文档真相源与向量派生索引的职责。
  \item \textbf{Python与前端}：Python、FastAPI、Vue 3、Next.js、TypeScript、微信小程序；完成过 AI 应用前后端闭环。
  \item \textbf{模型基础}：理解 Transformer、Attention、Embedding 与编码器-解码器；完成 BERT/LSTM 中文 NER 和 BART 文本生成训练实验。
  \item \textbf{工程实践}：Docker、Nginx、Linux、Git、Micrometer/Actuator、Mock 与端到端测试、环境变量注入和 AES-GCM 密钥保护。
\end{itemize}

\logosection{\faWrench}{项目经历}

\datedline{\textbf{408Master：面向计算机考研 408 的 AI 学习平台}}{项目队长 · 校实训二等奖}
\datedline{\tripleInfo{Java 21 / Spring Boot 4}{Spring AI 2 / Qdrant / Redis}{MySQL / Vue 3 / 微信小程序}}{}
\datedline{\href{https://github.com/tiankai-333/master408}{github.com/tiankai-333/master408}}{}

\textbf{项目简介：}在开源考试系统基础上进行深度二次开发，覆盖 408 题库、考试、错题、知识图谱、个人画像、AI 解析、练习规划和组卷等功能。在完成产品闭环后，继续将旧 Java 8 / Spring Boot 2.1 项目中手写的模型调用、双轨 RAG 和规则编排，渐进升级为 Java 21、Spring Boot 4 与 Spring AI 2 架构；大量复用已有 Controller、DTO、Mapper、业务 Service、前端及题库资产。

\textbf{核心职责：}
\begin{itemize}[parsep=0.5ex]
  \item \textbf{模型与记忆}：以 Spring AI \texttt{ChatModel/ChatClient} 统一同步和 SSE 流式调用，并兼容 DeepSeek、智谱等 OpenAI 协议模型；使用 Redis ChatMemory 按认证用户和随机会话 ID 隔离，配置 12 条消息窗口、24 小时 TTL、新对话和本人记忆清除。
  \item \textbf{多轮教学}：重写矛盾的柏拉图式 Prompt，使模型每轮只提出一个关键问题并根据学生回答继续追问；真实 DeepSeek 三轮验证前两轮不泄露答案，学生明确要求后才讲解。
  \item \textbf{RAG 与数据}：以 MySQL 为文档真相源、Qdrant 为可重建派生索引，在线查询统一为 Spring AI \texttt{VectorStore}；删除手写 Qdrant HTTP、MySQL 全量向量加载和 JVM O(N) 余弦兜底。将超长资料规范化为 6818 个 chunk，当前业务 1428 个全部索引成功。
  \item \textbf{混合编排与真实工具}：显式 intent、高置信规则和强上下文走确定性路径，低置信文本由模型选择 \texttt{plan\_practice}；复用旧 Planner、题库 SQL 和事务建卷，记录真实工具参数、结果、耗时和错误。
  \item \textbf{副作用控制}：模型只能生成练习草案，用户确认后才能建卷；题数上限、用户身份、事务和业务校验仍由 Java Service 控制，工具对象按请求绑定当前用户。
\end{itemize}

\textbf{技术难点与解决方案：}

\textbf{难点 1：旧业务资产完整，同时重写框架、业务和 AI 层会使错误无法定位}

\textbf{解决方案：}采用渐进迁移，保留 API、Mapper、题库与成熟 Service；旧调用和 Spring AI 新调用先经同一接口切换，完成隔离与真实模型验证后再删除旧实现，避免为目录整齐重写可靠业务。

\vspace{0.5ex}
\textbf{难点 2：ChatMemory 可能发生多用户污染，且网页历史容易与 Prompt Memory 混淆}

\textbf{解决方案：}服务端用认证 userId 再次绑定会话 ID，客户端不能提交最终 Redis Key；新对话轮换随机 ID，清除接口只能删除本人记忆。Redis Memory 只服务于模型上下文，网页历史保持独立职责。

\vspace{0.5ex}
\textbf{难点 3：旧 RAG 元数据显示已索引，但新 Qdrant 实际为空，超长文本还会使 Embedding 失败}

\textbf{解决方案：}将 embedding 模型和 collection 纳入索引版本，以 \texttt{force + afterId} 游标重建；先检查数据分布，再规范化超长资料并批量写入。最终 Qdrant 与 MySQL 新版本 indexed 元数据均为 1428，使用四类真实 408 查询验证召回。

\vspace{0.5ex}
\textbf{难点 4：纯规则覆盖有限，完全交给模型又会增加费用、延迟和误执行风险}

\textbf{解决方案：}高置信请求走零额外调用的确定性 fast path，模糊表达才进入 Tool Calling；模型只提出窄 Schema 动作，Java 应用负责执行、权限、事务和审计，写操作必须人工确认。

\textbf{阶段成果：}
\begin{itemize}[parsep=0.5ex]
  \item RAG 当前业务 1428 个 chunk 全部进入 Qdrant，索引失败 0；四类真实 408 查询均召回相关结果。
  \item 验证“模型请求工具—应用执行—结果回传模型—审计日志”循环，后端自动测试 28 项全部通过。
\end{itemize}

\datedline{\textbf{StoryCraft：AI 文字冒险游戏（已上线）}}{}
\datedline{\tripleInfo{Python / FastAPI}{Next.js / TypeScript}{Agent Runtime / SSE}}{}
\datedline{\href{https://github.com/baodashuaige/StoryCraft}{github.com/baodashuaige/StoryCraft} \quad
\href{https://game.mastertooling.shop}{game.mastertooling.shop}}{}

\textbf{项目简介：}构建可在线游玩的 AI 文字冒险游戏。NPC 具有人格、秘密、信任值、线索和剧情记忆，玩家行为影响多结局；提供 Web 与命令行入口，并将 NPC 对话能力抽取为 npm 包 \texttt{@wutiankai/npc-dialogue}。

\textbf{核心职责：}
\begin{itemize}[parsep=0.5ex]
  \item 设计 Agent Runtime，划分“LLM 生成叙事与候选动作、运行时校验、确定性代码提交状态”的职责；以 Prompt 和门控限制感知范围与可用动作。
  \item 完成 REST、SSE、前端交互、服务端密钥加密、测试和部署；状态只在完整响应校验后提交，避免流式中断产生半完成存档，自动化测试 176 passed。
\end{itemize}

\textbf{技术难点与解决方案：}

\textbf{难点：}概率模型可能虚构物品、跳过前置条件或替玩家作出选择，叙事自由度与确定性游戏规则冲突。

\textbf{解决方案：}Prompt 只提供角色应知信息和合法动作候选，模型输出先经过 Agent Runtime 门控，非法动作不进入状态机；SSE 只展示候选叙事，状态在完整结果校验后原子提交，使连接失败时世界状态保持一致。

\textbf{项目成果：}产品已部署并可公网访问；NPC 对话模块形成独立 npm 包，通过实际游戏验证“模型生成与确定性执行分离”的 Agent 设计。

\datedline{\textbf{自然语言处理课程项目：中文 NER 与文本生成}}{课程项目}
\datedline{\tripleInfo{BERT / LSTM}{BART}{CLUENER2020}}{}
\begin{itemize}[parsep=0.5ex]
  \item 基于 CLUENER2020 完成中文实体标签处理，分别使用 LSTM 与 BERT 训练命名实体识别模型，对比序列建模与预训练上下文表征。
  \item 基于 BART 进行文本生成微调，完成数据处理、训练、参数调整、推理与结果评估，理解模型微调与 Prompt、RAG、Memory、Tool Calling分别改变哪一层能力。
\end{itemize}

\logosection{\faTrophy}{荣誉与证书}
\datedline{东华大学软件工程实训二等奖（408Master 项目队长）}{}
\datedline{人工智能训练师职业技能培训}{预计 2026.08 取证}
